\chapter{Blood in Urine}

\section*{What Is This Test?}
Testing for blood in urine, or hematuria evaluation, determines whether red blood cells or heme are present and investigates the urinary tract as a possible source. Visible blood and blood detected only by testing require clinical assessment.

\section*{Alternative Names}
\begin{itemize}
  \item Hematuria Testing
  \item Urine Blood Test
  \item Microscopic Hematuria Evaluation
\end{itemize}

\section*{Principle}
A urine dipstick detects heme pigment, while microscopic examination identifies and counts red blood cells and may show casts or dysmorphic cells. Urine culture, kidney-function tests, and imaging or cystoscopy may be added according to symptoms and risk.

\section*{Why This Test Is Ordered}
It is used for pink, red, or cola-colored urine; dipstick blood; urinary pain; flank pain; recurrent urinary infection; stones; trauma; or concern for kidney disease or urinary-tract cancer.

\section*{Primary vs Secondary Test}
Urinalysis with microscopy is the primary test. A positive dipstick should generally be confirmed with microscopy because hemoglobin, myoglobin, menstrual contamination, concentrated urine, and some foods or medicines can affect color or results.

\section*{How This Test Is Performed}
A clean-catch midstream urine specimen is collected in a sterile container and examined promptly. The laboratory may perform dipstick testing, sediment microscopy, culture, and selected chemical measurements.

\section*{What Preparation Is Needed}
Follow clean-catch instructions. Tell the clinician about menstruation, vigorous exercise, recent trauma or procedures, anticoagulants, urinary symptoms, and medicines. Do not stop prescribed anticoagulants without medical advice.

\section*{Understanding Results}
Blood may arise from infection, stones, kidney inflammation, prostate disease, trauma, exercise, menstruation, or malignancy. Protein, casts, dysmorphic red cells, or reduced kidney function suggest a possible renal source. Persistent unexplained hematuria warrants risk-based evaluation.

\section*{Follow-up Tests}
Follow-up may include repeat urinalysis, urine culture, urine protein quantification, creatinine/eGFR, renal ultrasound or CT urography, urology evaluation, cystoscopy, or nephrology referral.

\section*{References}
\begin{enumerate}
  \item MedlinePlus. Blood in Urine.
  \item American Urological Association. Microhematuria guideline.
  \item National Kidney Foundation. Urinalysis and kidney disease information.
\end{enumerate}

\clearpage
\section*{Understanding Results and Follow-up}
The laboratory report should identify the specimen, method, units, reference interval or decision limit, and whether the result is positive, negative, abnormal, or inconclusive. Reference intervals vary with age, sex, pregnancy, specimen type, assay, and laboratory; a result outside a range is not automatically a diagnosis, and a result within a range does not always exclude disease. Discuss unexpected results with the ordering clinician, who can decide whether repeat or confirmatory testing, treatment, or referral is appropriate. Do not change medicines or delay urgent care based on an isolated result.


\section*{Regulatory Status and Authorization Date}
This chapter describes a test category, not a specific commercial product. FDA, CDSCO, EU IVDR, and MHRA approval or registration dates therefore cannot be assigned without the assay manufacturer, product name, instrument, and intended use. For laboratory-developed tests, an individual agency approval date may not exist. See the Preface for the interpretation of regulatory status.

\clearpage
